\scanpage{115}\mspage{58}

On voit donc que pour $\forall\,s\in S$, l'ordre du \deleted{\unsure{$s$}} sommet $s$ pour le graphe $A_\lambda$ est \deleted{\unsure{aussi}} égal à son ordre \inserted{$n$} pour $A_0$. Notons aussi que par transport des structures, $G_0^+$ \deleted{opère} transforme $A_\lambda$ en lui-=, i.e. $G_0^+$ opère par autom. du graphe $\Gamma_\lambda$. \deleted{ainsi} \deleted{qui est donc} On voit aussitôt qu'il est \inserted{simplement} transitif sur l'ens. des arêtes orientées de \deleted{\unsure{$G$}}$\Gamma_\lambda$, puisqu'il en est ainsi pour $\Gamma_0$.

On veut exprimer que $\Gamma_\lambda$ définit une structure de polyèdre régulier triangulé isomorphe à $P_0$. On va supposer pour ceci, d'abord, que \deleted{\unsure{le}}la structure de graphe $\Gamma_0\inserted{=(S_0,A_0)}$ détermine la structure de polyèdre triangulé $P_0=(S_0,F_0)$; \deleted{\unsure{p.ex.|i.e.}}

\[
(T_0\,4)\quad \deleted{\unsure{\text{on}}}\ \text{a}\quad F_0=\left\{f\in\unsure{\mathbb P_3|\mathcal P_3|\mathfrak P_3}(S)\ \middle|\ 
 \begin{gathered}\forall a\in\mathbb P_2(f),\ \text{on a}\\a\in A_0\end{gathered}\right\}
\tag*{$(38)$}
\]

% slanted
\margin{devrait venir avant $T_0\,1,T_0\,2,T_0\,3$ $\longleftarrow$}

(la structure de graphe détermine la structure triangulée)

J'ignore, s'il existe une structure de 2-polyèdre régulier \deleted{orientable} triangulé orientable, qui ne
% archive sheet number: 114
